Despite recent advances, large language models (LLMs) exhibit uneven performance across languages and often struggle with code-switching—the natural alternation between multiple languages. While code-switching is a common and meaningful mode of human communication, LLMs remain brittle to code-switching inputs, treating them as noise. This talk reframes code-switching as a valuable linguistic signal, exploring how it can be leveraged across the LLM lifecycle, from evaluation, to training, to inference. I will begin with output language alignment (OLA), highlighting systematic failures to respond in the expected language in code-switching contexts. I will also introduce code-switching red-teaming (CSRT), which reveals vulnerabilities in safety and multilingual understanding. I will then present code-switching curriculum learning (CSCL), a data-efficient approach to cross-lingual transfer inspired by human second-language learning. Finally, I will introduce code-switching in-context learning (CSICL), which uses gradual transitions from a target language to English to improve multilingual generalization without additional training.